%!TEX program = pdflatex
% IEEE Conference Paper Template
% 适用于IEEE会议论文

\documentclass[conference]{IEEEtran}

% 编码设置
\usepackage[utf8]{inputenc}
\usepackage[T1]{fontenc}

% 数学支持
\usepackage{amsmath}
\usepackage{amssymb}
\usepackage{amsfonts}
\usepackage{etoolbox}

% 图表支持
\usepackage{graphicx}
\usepackage{subfig}
\usepackage{float}
\usepackage{booktabs}

% 表格
\usepackage{tabularx}
\usepackage{colortbl}

% 代码环境
\usepackage{listings}

% 参考文献
\usepackage{cite}

% 超链接
\usepackage[colorlinks=true, linkcolor=black, citecolor=black, urlcolor=black]{hyperref}

% PDF元数据
\hypersetup{
    pdftitle={Paper Title},
    pdfauthor={Author Name},
    pdfkeywords={keyword1, keyword2, keyword3},
}

% 自定义命令
\DeclareMathOperator*{\argmax}{argmax}
\DeclareMathOperator*{\argmin}{argmin}
\providecommand{\IEEEauthorrefmark}[1]{$^{\ #1}$}

% Title
\title{Your Paper Title Here}

\author{
    \IEEEauthorblockN{Author Name}
    \IEEEauthorblockA{
        Department of Computer Science\\
        University Name\\
        City, Country\\
        email@example.com
    }
    \and
    \IEEEauthorblockN{Co-Author Name}
    \IEEEauthorblockA{
        Another Institution\\
        City, Country\\
        coauthor@example.com
    }
}

\begin{document}

\maketitle

\begin{abstract}
Your abstract goes here. It should summarize the main contributions of your paper in a concise manner. IEEE abstracts are typically limited to 150 words.
\end{abstract}

\begin{IEEEkeywords}
keyword1, keyword2, keyword3, keyword4
\end{IEEEkeywords}

\section{Introduction}
\label{sec:introduction}

Introduce the problem and motivation here. Explain why this work is important and what gap it fills in the literature.

\section{Related Work}
\label{sec:related}

Discuss related work here. Compare and contrast your approach with existing methods.

\section{Methodology}
\label{sec:method}

Explain your methodology in detail.

\subsection{Background}

Provide necessary background information.

\subsection{Proposed Approach}

Describe your proposed approach in detail.

\section{Experiments}
\label{sec:experiments}

Describe your experimental setup and results.

\subsection{Datasets}

Describe the datasets used in your experiments.

\subsection{Results}

Present your experimental results.

\begin{table}[htbp]
    \centering
    \caption{Example table}
    \label{tab:example}
    \begin{tabular}{lcc}
        \toprule
        Method & Accuracy & F1-Score \\
        \midrule
        Baseline 1 & 85.2 & 0.85 \\
        Baseline 2 & 87.1 & 0.87 \\
        Ours & \textbf{89.5} & \textbf{0.89} \\
        \bottomrule
    \end{tabular}
\end{table}

\begin{figure}[htbp]
    \centering
    \includegraphics[width=0.8\linewidth]{figures/fig1.png}
    \caption{Example figure}
    \label{fig:example}
\end{figure}

\section{Conclusion}
\label{sec:conclusion}

Summarize your work and discuss future directions.

\section*{Acknowledgment}

Thank collaborators and funding agencies.

% References
\begin{thebibliography}{1}

    \bibitem{ref1}
    Author, ``Title of the paper,'' in \emph{Proceedings of the Conference Name}, Year, pp. XX--XX.

    \bibitem{ref2}
    Author and Co-Author, ``Title of another paper,'' \emph{Journal Name}, vol. X, no. X, pp. XX--XX, Year.

\end{thebibliography}

\end{document}
